%%%%%%%%%%%%%%%%%%%%%%%%%%%%%%%%%%%%%%%%%%%%%%%%%%%%%%%%%%%%
% 11.1 NUMERICAL ANALYSIS
% The Composition of Advanced Mathematics
%%%%%%%%%%%%%%%%%%%%%%%%%%%%%%%%%%%%%%%%%%%%%%%%%%%%%%%%%%%%

\section{Numerical Analysis}
\label{sec:numerical-analysis}

\prerequisite{
Calculus, linear algebra, sequences and series, differential equations,
Taylor expansions, optimization fundamentals, and basic programming
concepts.
}

\learningobjectives{
By the end of this section, the reader should be able to:
\begin{itemize}
    \item explain the purpose of numerical analysis;
    \item distinguish exact and approximate computation;
    \item identify sources of numerical error;
    \item distinguish absolute and relative error;
    \item understand significant digits and machine precision;
    \item understand floating-point representation;
    \item recognize overflow, underflow, and rounding error;
    \item distinguish truncation error from roundoff error;
    \item understand catastrophic cancellation;
    \item define conditioning of a mathematical problem;
    \item distinguish conditioning from numerical stability;
    \item understand forward and backward error;
    \item use Taylor expansions to derive numerical approximations;
    \item define consistency, stability, and convergence;
    \item understand convergence order;
    \item solve nonlinear scalar equations numerically;
    \item apply the bisection method;
    \item apply fixed-point iteration;
    \item apply Newton's method;
    \item apply the secant method;
    \item compare bracketing and open root-finding methods;
    \item understand stopping criteria;
    \item approximate derivatives using finite differences;
    \item derive forward, backward, and centered differences;
    \item understand finite-difference error order;
    \item apply Richardson extrapolation;
    \item approximate definite integrals numerically;
    \item apply midpoint and trapezoidal rules;
    \item apply Simpson's rule;
    \item understand composite quadrature;
    \item estimate numerical-integration error;
    \item understand adaptive quadrature conceptually;
    \item distinguish interpolation from approximation;
    \item construct polynomial interpolants;
    \item understand Lagrange interpolation;
    \item understand Newton divided differences;
    \item state the polynomial interpolation error formula;
    \item recognize Runge's phenomenon;
    \item understand piecewise polynomial interpolation;
    \item understand cubic splines conceptually;
    \item understand least-squares approximation;
    \item distinguish interpolation from least-squares fitting;
    \item recognize numerical conditioning in interpolation;
    \item understand iterative refinement conceptually;
    \item recognize residual-based error checks;
    \item understand computational complexity in numerical methods;
    \item understand verification and validation of numerical results;
    \item connect numerical analysis with scientific computing,
          differential equations, optimization, simulation, statistics,
          and applied mathematics.
\end{itemize}
}

%%%%%%%%%%%%%%%%%%%%%%%%%%%%%%%%%%%%%%%%%%%%%%%%%%%%%%%%%%%%
\subsection{What Is Numerical Analysis?}
\label{subsec:what-is-numerical-analysis}
%%%%%%%%%%%%%%%%%%%%%%%%%%%%%%%%%%%%%%%%%%%%%%%%%%%%%%%%%%%%

Numerical analysis studies algorithms for approximating mathematical
quantities that cannot be computed exactly or conveniently.

Examples include:

\[
\boxed{
\text{roots of nonlinear equations},
}
\]

\[
\boxed{
\text{solutions of differential equations},
}
\]

\[
\boxed{
\text{large systems of equations},
}
\]

\[
\boxed{
\text{definite integrals},
}
\]

and

\[
\boxed{
\text{optimization problems}.
}
\]

The central questions are not only:

\[
\boxed{
\text{Does the algorithm produce an answer?}
}
\]

but also:

\[
\boxed{
\text{How accurate is it?}
}
\]

\[
\boxed{
\text{How stable is it?}
}
\]

and

\[
\boxed{
\text{How much computation does it require?}
}
\]

%%%%%%%%%%%%%%%%%%%%%%%%%%%%%%%%%%%%%%%%%%%%%%%%%%%%%%%%%%%%
\subsection{Exact Versus Numerical Solutions}
\label{subsec:exact-versus-numerical}
%%%%%%%%%%%%%%%%%%%%%%%%%%%%%%%%%%%%%%%%%%%%%%%%%%%%%%%%%%%%

An exact solution may be represented symbolically.

For example,

\[
x^2-2=0
\]

has exact positive solution

\[
\boxed{
x=\sqrt{2}.
}
\]

A numerical approximation might be

\[
\boxed{
x\approx1.41421356.
}
\]

Numerical methods replace exact symbolic expressions with approximations
that can be computed to a desired accuracy.

%%%%%%%%%%%%%%%%%%%%%%%%%%%%%%%%%%%%%%%%%%%%%%%%%%%%%%%%%%%%
\subsection{Sources of Numerical Error}
\label{subsec:sources-numerical-error}
%%%%%%%%%%%%%%%%%%%%%%%%%%%%%%%%%%%%%%%%%%%%%%%%%%%%%%%%%%%%

Numerical error can arise from several sources:

\[
\boxed{
\text{data error},
}
\]

\[
\boxed{
\text{modeling error},
}
\]

\[
\boxed{
\text{truncation error},
}
\]

\[
\boxed{
\text{roundoff error},
}
\]

and

\[
\boxed{
\text{iteration error}.
}
\]

These errors should not be treated as interchangeable.

%%%%%%%%%%%%%%%%%%%%%%%%%%%%%%%%%%%%%%%%%%%%%%%%%%%%%%%%%%%%
\subsection{Absolute Error}
\label{subsec:absolute-error-numerical}
%%%%%%%%%%%%%%%%%%%%%%%%%%%%%%%%%%%%%%%%%%%%%%%%%%%%%%%%%%%%

If \(x\) is the exact value and \(\widehat{x}\) is an approximation, the
absolute error is

\[
\boxed{
E_{\mathrm{abs}}
=
|x-\widehat{x}|.
}
\]

Absolute error measures the size of the error in the units of the
quantity.

%%%%%%%%%%%%%%%%%%%%%%%%%%%%%%%%%%%%%%%%%%%%%%%%%%%%%%%%%%%%
\subsection{Relative Error}
\label{subsec:relative-error-numerical}
%%%%%%%%%%%%%%%%%%%%%%%%%%%%%%%%%%%%%%%%%%%%%%%%%%%%%%%%%%%%

For \(x\neq0\), the relative error is

\[
\boxed{
E_{\mathrm{rel}}
=
\frac{
|x-\widehat{x}|
}{
|x|
}.
}
\]

Relative error measures the error compared with the scale of the exact
quantity.

A percentage error is

\[
\boxed{
100E_{\mathrm{rel}}\%.
}
\]

%%%%%%%%%%%%%%%%%%%%%%%%%%%%%%%%%%%%%%%%%%%%%%%%%%%%%%%%%%%%
\subsection{Worked Example: Absolute and Relative Error}
\label{subsec:worked-absolute-relative-error}
%%%%%%%%%%%%%%%%%%%%%%%%%%%%%%%%%%%%%%%%%%%%%%%%%%%%%%%%%%%%

\begin{workedexamplebox}{Approximating \(\pi\)}

Suppose

\[
\widehat{\pi}=3.14.
\]

The absolute error is

\[
|\pi-3.14|
\approx
0.00159265.
\]

The relative error is

\[
\frac{
|\pi-3.14|
}{
\pi
}
\approx
0.00050696.
\]

Thus the percentage relative error is approximately

\[
0.0507\%.
\]

\begin{answerbox}
\[
\boxed{
E_{\mathrm{abs}}
\approx0.00159265
}
\]

and

\[
\boxed{
E_{\mathrm{rel}}
\approx5.07\times10^{-4}.
}
\]
\end{answerbox}

\end{workedexamplebox}

%%%%%%%%%%%%%%%%%%%%%%%%%%%%%%%%%%%%%%%%%%%%%%%%%%%%%%%%%%%%
\subsection{Significant Digits}
\label{subsec:significant-digits}
%%%%%%%%%%%%%%%%%%%%%%%%%%%%%%%%%%%%%%%%%%%%%%%%%%%%%%%%%%%%

An approximation has \(p\) correct significant decimal digits when its
relative error is sufficiently small compared with

\[
10^{-p}.
\]

In practice, significant digits describe how many leading digits of a
numerical value are reliable.

They are distinct from the number of digits merely displayed by software.

%%%%%%%%%%%%%%%%%%%%%%%%%%%%%%%%%%%%%%%%%%%%%%%%%%%%%%%%%%%%
\subsection{Floating-Point Representation}
\label{subsec:floating-point-representation}
%%%%%%%%%%%%%%%%%%%%%%%%%%%%%%%%%%%%%%%%%%%%%%%%%%%%%%%%%%%%

Computers store most real numbers approximately using floating-point
representation.

A simplified normalized form is

\[
\boxed{
x
=
\pm
m
\beta^e,
}
\]

where:

\[
m
=
\text{mantissa or significand},
\]

\[
\beta
=
\text{base},
\]

and

\[
e
=
\text{exponent}.
\]

In modern computers,

\[
\boxed{
\beta=2
}
\]

is standard.

%%%%%%%%%%%%%%%%%%%%%%%%%%%%%%%%%%%%%%%%%%%%%%%%%%%%%%%%%%%%
\subsection{Finite Precision}
\label{subsec:finite-precision}
%%%%%%%%%%%%%%%%%%%%%%%%%%%%%%%%%%%%%%%%%%%%%%%%%%%%%%%%%%%%

Only finitely many bits are available to represent a floating-point
number.

Therefore most real numbers cannot be represented exactly.

For example, just as

\[
\frac13
=
0.3333\ldots
\]

has no finite decimal expansion,

many simple decimal fractions have no finite binary representation.

%%%%%%%%%%%%%%%%%%%%%%%%%%%%%%%%%%%%%%%%%%%%%%%%%%%%%%%%%%%%
\subsection{Rounding}
\label{subsec:floating-point-rounding}
%%%%%%%%%%%%%%%%%%%%%%%%%%%%%%%%%%%%%%%%%%%%%%%%%%%%%%%%%%%%

Let

\[
\operatorname{fl}(x)
\]

denote the floating-point representation of \(x\).

A common model is

\[
\boxed{
\operatorname{fl}(x)
=
x(1+\delta),
}
\]

where

\[
\boxed{
|\delta|
\leq
u.
}
\]

The quantity \(u\) is the unit roundoff under the chosen floating-point
system and rounding convention.

%%%%%%%%%%%%%%%%%%%%%%%%%%%%%%%%%%%%%%%%%%%%%%%%%%%%%%%%%%%%
\subsection{Machine Precision}
\label{subsec:machine-precision}
%%%%%%%%%%%%%%%%%%%%%%%%%%%%%%%%%%%%%%%%%%%%%%%%%%%%%%%%%%%%

Informally, machine precision describes the scale at which floating-point
rounding prevents neighboring represented values near \(1\) from being
distinguished.

Terminology varies between sources.

A common computational quantity is machine epsilon:

\[
\boxed{
\varepsilon_{\mathrm{mach}}
=
\text{distance from }1
\text{ to the next larger representable number}.
}
\]

For IEEE double precision, this value is approximately

\[
\boxed{
2.22\times10^{-16}.
}
\]

%%%%%%%%%%%%%%%%%%%%%%%%%%%%%%%%%%%%%%%%%%%%%%%%%%%%%%%%%%%%
\subsection{Roundoff Error}
\label{subsec:roundoff-error}
%%%%%%%%%%%%%%%%%%%%%%%%%%%%%%%%%%%%%%%%%%%%%%%%%%%%%%%%%%%%

Roundoff error occurs because arithmetic is performed with finite
precision.

Even if the mathematical operation is exact,

\[
x+y,
\]

the computer actually stores an approximation such as

\[
\boxed{
\operatorname{fl}
(
\operatorname{fl}(x)
+
\operatorname{fl}(y)
).
}
\]

Repeated arithmetic operations can accumulate rounding effects.

%%%%%%%%%%%%%%%%%%%%%%%%%%%%%%%%%%%%%%%%%%%%%%%%%%%%%%%%%%%%
\subsection{Overflow}
\label{subsec:numerical-overflow}
%%%%%%%%%%%%%%%%%%%%%%%%%%%%%%%%%%%%%%%%%%%%%%%%%%%%%%%%%%%%

Overflow occurs when a computed magnitude exceeds the largest
representable floating-point number.

Conceptually,

\[
\boxed{
|x|
>
x_{\max}.
}
\]

Depending on the system, the result may become an infinity value or raise
an error.

%%%%%%%%%%%%%%%%%%%%%%%%%%%%%%%%%%%%%%%%%%%%%%%%%%%%%%%%%%%%
\subsection{Underflow}
\label{subsec:numerical-underflow}
%%%%%%%%%%%%%%%%%%%%%%%%%%%%%%%%%%%%%%%%%%%%%%%%%%%%%%%%%%%%

Underflow occurs when a nonzero magnitude is too small for ordinary
normalized representation.

Modern floating-point systems use subnormal numbers to represent some
values below the normal range.

Extremely small values may eventually round to zero.

%%%%%%%%%%%%%%%%%%%%%%%%%%%%%%%%%%%%%%%%%%%%%%%%%%%%%%%%%%%%
\subsection{Catastrophic Cancellation}
\label{subsec:catastrophic-cancellation}
%%%%%%%%%%%%%%%%%%%%%%%%%%%%%%%%%%%%%%%%%%%%%%%%%%%%%%%%%%%%

Catastrophic cancellation occurs when nearly equal floating-point numbers
are subtracted.

Suppose

\[
x
\approx
y.
\]

Then

\[
x-y
\]

may be much smaller than either input.

Errors already present in the leading digits can then dominate the
remaining result.

%%%%%%%%%%%%%%%%%%%%%%%%%%%%%%%%%%%%%%%%%%%%%%%%%%%%%%%%%%%%
\subsection{Worked Example: Stable Algebraic Reformulation}
\label{subsec:worked-stable-reformulation}
%%%%%%%%%%%%%%%%%%%%%%%%%%%%%%%%%%%%%%%%%%%%%%%%%%%%%%%%%%%%

\begin{workedexamplebox}{Avoiding Cancellation}

Consider

\[
f(x)
=
\sqrt{x+1}-\sqrt{x}
\]

for large \(x\).

Direct subtraction involves two nearly equal quantities.

Rationalize:

\[
f(x)
=
\frac{
(\sqrt{x+1}-\sqrt{x})
(\sqrt{x+1}+\sqrt{x})
}{
\sqrt{x+1}+\sqrt{x}
}.
\]

Therefore,

\[
f(x)
=
\frac{
1
}{
\sqrt{x+1}+\sqrt{x}
}.
\]

\begin{answerbox}
\[
\boxed{
\sqrt{x+1}-\sqrt{x}
=
\frac{
1
}{
\sqrt{x+1}+\sqrt{x}
}
}
\]

and the second expression is usually numerically preferable for large
\(x\).
\end{answerbox}

\end{workedexamplebox}

%%%%%%%%%%%%%%%%%%%%%%%%%%%%%%%%%%%%%%%%%%%%%%%%%%%%%%%%%%%%
\subsection{Truncation Error}
\label{subsec:truncation-error}
%%%%%%%%%%%%%%%%%%%%%%%%%%%%%%%%%%%%%%%%%%%%%%%%%%%%%%%%%%%%

Truncation error occurs when an infinite mathematical process is replaced
by a finite approximation.

For example,

\[
e^x
=
\sum_{k=0}^{\infty}
\frac{x^k}{k!}.
\]

Approximating with

\[
\boxed{
e^x
\approx
\sum_{k=0}^{n}
\frac{x^k}{k!}
}
\]

creates truncation error.

%%%%%%%%%%%%%%%%%%%%%%%%%%%%%%%%%%%%%%%%%%%%%%%%%%%%%%%%%%%%
\subsection{Taylor's Theorem}
\label{subsec:taylor-theorem-numerical}
%%%%%%%%%%%%%%%%%%%%%%%%%%%%%%%%%%%%%%%%%%%%%%%%%%%%%%%%%%%%

For sufficiently differentiable \(f\),

\[
f(x+h)
=
f(x)
+
hf'(x)
+
\frac{h^2}{2!}f''(x)
+
\cdots
+
\frac{h^n}{n!}f^{(n)}(x)
+
R_{n+1}.
\]

A common remainder form is

\[
\boxed{
R_{n+1}
=
\frac{
f^{(n+1)}(\xi)
}{
(n+1)!
}
h^{n+1}
}
\]

for some \(\xi\) between \(x\) and \(x+h\).

Taylor expansions are one of the main tools for deriving and analyzing
numerical methods.

%%%%%%%%%%%%%%%%%%%%%%%%%%%%%%%%%%%%%%%%%%%%%%%%%%%%%%%%%%%%
\subsection{Big-O Error Notation}
\label{subsec:big-o-numerical-error}
%%%%%%%%%%%%%%%%%%%%%%%%%%%%%%%%%%%%%%%%%%%%%%%%%%%%%%%%%%%%

If an error satisfies

\[
E(h)
=
O(h^p)
\]

as \(h\to0\), then there are constants \(C\) and \(h_0>0\) such that

\[
\boxed{
|E(h)|
\leq
C|h|^p
}
\]

for sufficiently small \(h\).

The exponent \(p\) describes the asymptotic error order.

%%%%%%%%%%%%%%%%%%%%%%%%%%%%%%%%%%%%%%%%%%%%%%%%%%%%%%%%%%%%
\subsection{Convergence Order}
\label{subsec:convergence-order}
%%%%%%%%%%%%%%%%%%%%%%%%%%%%%%%%%%%%%%%%%%%%%%%%%%%%%%%%%%%%

Suppose approximations \(x_n\) converge to \(x^*\).

If

\[
\boxed{
|x_{n+1}-x^*|
\approx
C
|x_n-x^*|^p,
}
\]

then \(p\) is the order of convergence.

Important cases include:

\[
\boxed{
p=1
\quad
\text{linear convergence},
}
\]

and

\[
\boxed{
p=2
\quad
\text{quadratic convergence}.
}
\]

%%%%%%%%%%%%%%%%%%%%%%%%%%%%%%%%%%%%%%%%%%%%%%%%%%%%%%%%%%%%
\subsection{Conditioning}
\label{subsec:numerical-conditioning}
%%%%%%%%%%%%%%%%%%%%%%%%%%%%%%%%%%%%%%%%%%%%%%%%%%%%%%%%%%%%

Conditioning describes sensitivity of the exact mathematical solution to
small perturbations in the data.

A problem is well conditioned when small input changes produce small
output changes.

A problem is ill conditioned when small input changes can produce large
output changes.

Conditioning is a property of the problem, not the algorithm.

%%%%%%%%%%%%%%%%%%%%%%%%%%%%%%%%%%%%%%%%%%%%%%%%%%%%%%%%%%%%
\subsection{Relative Condition Number}
\label{subsec:relative-condition-number}
%%%%%%%%%%%%%%%%%%%%%%%%%%%%%%%%%%%%%%%%%%%%%%%%%%%%%%%%%%%%

For a differentiable scalar function

\[
y=f(x),
\]

a relative condition number can be defined as

\[
\boxed{
\kappa_f(x)
=
\left|
\frac{
xf'(x)
}{
f(x)
}
\right|
}
\]

when the expression is defined.

Approximately,

\[
\boxed{
\frac{|\Delta y|}{|y|}
\approx
\kappa_f(x)
\frac{|\Delta x|}{|x|}.
}
\]

%%%%%%%%%%%%%%%%%%%%%%%%%%%%%%%%%%%%%%%%%%%%%%%%%%%%%%%%%%%%
\subsection{Worked Example: Conditioning}
\label{subsec:worked-conditioning}
%%%%%%%%%%%%%%%%%%%%%%%%%%%%%%%%%%%%%%%%%%%%%%%%%%%%%%%%%%%%

\begin{workedexamplebox}{Conditioning of \(f(x)=x^n\)}

Let

\[
f(x)=x^n.
\]

Then

\[
f'(x)=nx^{n-1}.
\]

Hence,

\[
\kappa_f(x)
=
\left|
\frac{
x(nx^{n-1})
}{
x^n
}
\right|
=
n.
\]

Therefore,

\begin{answerbox}
\[
\boxed{
\kappa_f=n.
}
\]
\end{answerbox}

A small relative perturbation in \(x\) is amplified approximately by a
factor of \(n\) in the output.

\end{workedexamplebox}

%%%%%%%%%%%%%%%%%%%%%%%%%%%%%%%%%%%%%%%%%%%%%%%%%%%%%%%%%%%%
\subsection{Numerical Stability}
\label{subsec:numerical-stability}
%%%%%%%%%%%%%%%%%%%%%%%%%%%%%%%%%%%%%%%%%%%%%%%%%%%%%%%%%%%%

Stability describes how an algorithm propagates computational errors.

A numerically stable algorithm does not amplify errors unnecessarily.

Thus:

\[
\boxed{
\text{conditioning}
=
\text{property of the problem},
}
\]

while

\[
\boxed{
\text{stability}
=
\text{property of the algorithm}.
}
\]

%%%%%%%%%%%%%%%%%%%%%%%%%%%%%%%%%%%%%%%%%%%%%%%%%%%%%%%%%%%%
\subsection{Forward Error}
\label{subsec:forward-error}
%%%%%%%%%%%%%%%%%%%%%%%%%%%%%%%%%%%%%%%%%%%%%%%%%%%%%%%%%%%%

Forward error compares the computed solution

\[
\widehat{x}
\]

with the exact solution

\[
x^*.
\]

For example,

\[
\boxed{
E_{\mathrm{forward}}
=
|\widehat{x}-x^*|.
}
\]

This directly measures output accuracy.

%%%%%%%%%%%%%%%%%%%%%%%%%%%%%%%%%%%%%%%%%%%%%%%%%%%%%%%%%%%%
\subsection{Backward Error}
\label{subsec:backward-error}
%%%%%%%%%%%%%%%%%%%%%%%%%%%%%%%%%%%%%%%%%%%%%%%%%%%%%%%%%%%%

Backward error asks:

\[
\boxed{
\text{For what nearby input data is the computed answer exact?}
}
\]

An algorithm is backward stable if its computed result is the exact
solution of a problem with only a small perturbation of the original data.

Backward stability is a powerful numerical-analysis concept.

%%%%%%%%%%%%%%%%%%%%%%%%%%%%%%%%%%%%%%%%%%%%%%%%%%%%%%%%%%%%
\subsection{Residual}
\label{subsec:numerical-residual}
%%%%%%%%%%%%%%%%%%%%%%%%%%%%%%%%%%%%%%%%%%%%%%%%%%%%%%%%%%%%

For an equation

\[
f(x)=0,
\]

the residual of approximation \(\widehat{x}\) is

\[
\boxed{
r
=
f(\widehat{x}).
}
\]

A small residual means the approximation nearly satisfies the equation.

However, a small residual does not always imply small solution error if the
problem is ill conditioned.

%%%%%%%%%%%%%%%%%%%%%%%%%%%%%%%%%%%%%%%%%%%%%%%%%%%%%%%%%%%%
\subsection{Root-Finding Problems}
\label{subsec:root-finding-problems}
%%%%%%%%%%%%%%%%%%%%%%%%%%%%%%%%%%%%%%%%%%%%%%%%%%%%%%%%%%%%

A nonlinear root-finding problem seeks

\[
\boxed{
x^*
}
\]

such that

\[
\boxed{
f(x^*)=0.
}
\]

Many nonlinear equations cannot be solved symbolically.

Numerical root-finding methods generate a sequence

\[
x_0,x_1,x_2,\ldots
\]

intended to converge to a root.

%%%%%%%%%%%%%%%%%%%%%%%%%%%%%%%%%%%%%%%%%%%%%%%%%%%%%%%%%%%%
\subsection{Bisection Method}
\label{subsec:bisection-method}
%%%%%%%%%%%%%%%%%%%%%%%%%%%%%%%%%%%%%%%%%%%%%%%%%%%%%%%%%%%%

Suppose \(f\) is continuous on

\[
[a,b]
\]

and

\[
\boxed{
f(a)f(b)<0.
}
\]

The Intermediate Value Theorem guarantees at least one root in

\[
(a,b).
\]

Bisection repeatedly halves the interval.

%%%%%%%%%%%%%%%%%%%%%%%%%%%%%%%%%%%%%%%%%%%%%%%%%%%%%%%%%%%%
\subsection{Bisection Algorithm}
\label{subsec:bisection-algorithm}
%%%%%%%%%%%%%%%%%%%%%%%%%%%%%%%%%%%%%%%%%%%%%%%%%%%%%%%%%%%%

At each iteration calculate

\[
\boxed{
c
=
\frac{a+b}{2}.
}
\]

If

\[
f(c)=0,
\]

stop.

Otherwise retain whichever half interval contains a sign change.

Thus the bracket width is halved at every iteration.

%%%%%%%%%%%%%%%%%%%%%%%%%%%%%%%%%%%%%%%%%%%%%%%%%%%%%%%%%%%%
\subsection{Bisection Error Bound}
\label{subsec:bisection-error}
%%%%%%%%%%%%%%%%%%%%%%%%%%%%%%%%%%%%%%%%%%%%%%%%%%%%%%%%%%%%

If the initial interval has width

\[
b-a,
\]

then after \(n\) bisections the interval width is

\[
\boxed{
\frac{
b-a
}{
2^n
}.
}
\]

If the midpoint is used as the approximation, a corresponding error bound
is

\[
\boxed{
|x_n-x^*|
\leq
\frac{
b-a
}{
2^{n+1}
}
}
\]

under the usual indexing convention.

%%%%%%%%%%%%%%%%%%%%%%%%%%%%%%%%%%%%%%%%%%%%%%%%%%%%%%%%%%%%
\subsection{Worked Example: Bisection}
\label{subsec:worked-bisection}
%%%%%%%%%%%%%%%%%%%%%%%%%%%%%%%%%%%%%%%%%%%%%%%%%%%%%%%%%%%%

\begin{workedexamplebox}{Root of \(x^2-2\)}

Let

\[
f(x)=x^2-2.
\]

On

\[
[1,2],
\]

\[
f(1)=-1,
\qquad
f(2)=2.
\]

The first midpoint is

\[
c_1=1.5.
\]

Since

\[
f(1.5)=0.25>0,
\]

the root lies in

\[
[1,1.5].
\]

The next midpoint is

\[
c_2=1.25.
\]

Since

\[
f(1.25)<0,
\]

the root lies in

\[
[1.25,1.5].
\]

Continuing this procedure converges to

\[
\boxed{
\sqrt{2}.
}
\]

\end{workedexamplebox}

%%%%%%%%%%%%%%%%%%%%%%%%%%%%%%%%%%%%%%%%%%%%%%%%%%%%%%%%%%%%
\subsection{Advantages of Bisection}
\label{subsec:bisection-advantages}
%%%%%%%%%%%%%%%%%%%%%%%%%%%%%%%%%%%%%%%%%%%%%%%%%%%%%%%%%%%%

Bisection is attractive because it is:

\[
\boxed{
\text{simple},
}
\]

\[
\boxed{
\text{robust},
}
\]

and

\[
\boxed{
\text{guaranteed to converge}
}
\]

when a continuous sign-changing bracket is maintained.

Its main disadvantage is relatively slow convergence.

%%%%%%%%%%%%%%%%%%%%%%%%%%%%%%%%%%%%%%%%%%%%%%%%%%%%%%%%%%%%
\subsection{Fixed-Point Iteration}
\label{subsec:fixed-point-iteration}
%%%%%%%%%%%%%%%%%%%%%%%%%%%%%%%%%%%%%%%%%%%%%%%%%%%%%%%%%%%%

Rewrite

\[
f(x)=0
\]

as

\[
\boxed{
x=g(x).
}
\]

Then iterate

\[
\boxed{
x_{n+1}
=
g(x_n).
}
\]

A solution satisfying

\[
x^*=g(x^*)
\]

is called a fixed point.

%%%%%%%%%%%%%%%%%%%%%%%%%%%%%%%%%%%%%%%%%%%%%%%%%%%%%%%%%%%%
\subsection{Local Convergence of Fixed-Point Iteration}
\label{subsec:fixed-point-local-convergence}
%%%%%%%%%%%%%%%%%%%%%%%%%%%%%%%%%%%%%%%%%%%%%%%%%%%%%%%%%%%%

If \(g\) is differentiable near \(x^*\) and

\[
\boxed{
|g'(x^*)|<1,
}
\]

then fixed-point iteration is locally contractive and typically converges
for sufficiently close initial guesses.

If

\[
|g'(x^*)|>1,
\]

the fixed point is locally repelling.

%%%%%%%%%%%%%%%%%%%%%%%%%%%%%%%%%%%%%%%%%%%%%%%%%%%%%%%%%%%%
\subsection{Newton's Method}
\label{subsec:newtons-method}
%%%%%%%%%%%%%%%%%%%%%%%%%%%%%%%%%%%%%%%%%%%%%%%%%%%%%%%%%%%%

Newton's method approximates \(f\) near \(x_n\) by its tangent line:

\[
f(x)
\approx
f(x_n)
+
f'(x_n)(x-x_n).
\]

Setting the approximation equal to zero gives

\[
\boxed{
x_{n+1}
=
x_n
-
\frac{
f(x_n)
}{
f'(x_n)
}.
}
\]

%%%%%%%%%%%%%%%%%%%%%%%%%%%%%%%%%%%%%%%%%%%%%%%%%%%%%%%%%%%%
\subsection{Geometric Interpretation of Newton's Method}
\label{subsec:newton-geometric}
%%%%%%%%%%%%%%%%%%%%%%%%%%%%%%%%%%%%%%%%%%%%%%%%%%%%%%%%%%%%

At \(x_n\):

\begin{enumerate}
    \item draw the tangent line to \(y=f(x)\);
    \item find where the tangent intersects the \(x\)-axis;
    \item use that intercept as \(x_{n+1}\).
\end{enumerate}

Thus Newton's method repeatedly replaces the nonlinear equation with a
local linear approximation.

%%%%%%%%%%%%%%%%%%%%%%%%%%%%%%%%%%%%%%%%%%%%%%%%%%%%%%%%%%%%
\subsection{Worked Example: Newton's Method}
\label{subsec:worked-newton}
%%%%%%%%%%%%%%%%%%%%%%%%%%%%%%%%%%%%%%%%%%%%%%%%%%%%%%%%%%%%

\begin{workedexamplebox}{Approximating \(\sqrt{2}\)}

Let

\[
f(x)=x^2-2.
\]

Then

\[
f'(x)=2x.
\]

Newton's iteration is

\[
x_{n+1}
=
x_n
-
\frac{
x_n^2-2
}{
2x_n
}.
\]

Simplifying,

\[
\boxed{
x_{n+1}
=
\frac12
\left(
x_n+\frac{2}{x_n}
\right).
}
\]

Starting with

\[
x_0=1.5,
\]

we obtain

\[
x_1
=
\frac12
\left(
1.5+\frac{2}{1.5}
\right)
\approx1.4166667.
\]

Then

\[
x_2
\approx1.4142157.
\]

Then

\[
x_3
\approx1.41421356.
\]

\begin{answerbox}
\[
\boxed{
x\approx1.41421356.
}
\]
\end{answerbox}

\end{workedexamplebox}

%%%%%%%%%%%%%%%%%%%%%%%%%%%%%%%%%%%%%%%%%%%%%%%%%%%%%%%%%%%%
\subsection{Quadratic Convergence of Newton's Method}
\label{subsec:newton-quadratic-convergence}
%%%%%%%%%%%%%%%%%%%%%%%%%%%%%%%%%%%%%%%%%%%%%%%%%%%%%%%%%%%%

Near a simple root \(x^*\), under suitable smoothness assumptions,

\[
\boxed{
|e_{n+1}|
\approx
C|e_n|^2,
}
\]

where

\[
e_n=x_n-x^*.
\]

Thus Newton's method has quadratic local convergence.

Once sufficiently close to the root, the number of correct digits can
increase very rapidly.

%%%%%%%%%%%%%%%%%%%%%%%%%%%%%%%%%%%%%%%%%%%%%%%%%%%%%%%%%%%%
\subsection{Failure Modes of Newton's Method}
\label{subsec:newton-failure}
%%%%%%%%%%%%%%%%%%%%%%%%%%%%%%%%%%%%%%%%%%%%%%%%%%%%%%%%%%%%

Newton's method can fail when:

\[
\boxed{
f'(x_n)\approx0,
}
\]

\[
\boxed{
\text{the initial guess is poor},
}
\]

\[
\boxed{
\text{the iteration leaves the desired region},
}
\]

or

\[
\boxed{
\text{the function has complicated geometry}.
}
\]

Unlike bisection, Newton's method does not automatically preserve a root
bracket.

%%%%%%%%%%%%%%%%%%%%%%%%%%%%%%%%%%%%%%%%%%%%%%%%%%%%%%%%%%%%
\subsection{Multiple Roots}
\label{subsec:newton-multiple-roots}
%%%%%%%%%%%%%%%%%%%%%%%%%%%%%%%%%%%%%%%%%%%%%%%%%%%%%%%%%%%%

If \(x^*\) is a root of multiplicity \(m>1\), ordinary Newton iteration may
lose quadratic convergence.

A modified iteration is

\[
\boxed{
x_{n+1}
=
x_n
-
m
\frac{
f(x_n)
}{
f'(x_n)
}.
}
\]

When \(m\) is known, this can restore faster local convergence.

%%%%%%%%%%%%%%%%%%%%%%%%%%%%%%%%%%%%%%%%%%%%%%%%%%%%%%%%%%%%
\subsection{Secant Method}
\label{subsec:secant-method}
%%%%%%%%%%%%%%%%%%%%%%%%%%%%%%%%%%%%%%%%%%%%%%%%%%%%%%%%%%%%

The secant method approximates the derivative using two previous points:

\[
f'(x_n)
\approx
\frac{
f(x_n)-f(x_{n-1})
}{
x_n-x_{n-1}
}.
\]

Substituting into Newton's formula gives

\[
\boxed{
x_{n+1}
=
x_n
-
f(x_n)
\frac{
x_n-x_{n-1}
}{
f(x_n)-f(x_{n-1})
}.
}
\]

%%%%%%%%%%%%%%%%%%%%%%%%%%%%%%%%%%%%%%%%%%%%%%%%%%%%%%%%%%%%
\subsection{Advantages of the Secant Method}
\label{subsec:secant-advantages}
%%%%%%%%%%%%%%%%%%%%%%%%%%%%%%%%%%%%%%%%%%%%%%%%%%%%%%%%%%%%

The secant method does not require an explicit derivative.

It generally converges faster than bisection but less rapidly than
quadratically convergent Newton iteration.

Its local convergence order for a simple root is approximately

\[
\boxed{
\frac{
1+\sqrt{5}
}{
2
}
\approx1.618.
}
\]

%%%%%%%%%%%%%%%%%%%%%%%%%%%%%%%%%%%%%%%%%%%%%%%%%%%%%%%%%%%%
\subsection{Bracketing Versus Open Methods}
\label{subsec:bracketing-versus-open}
%%%%%%%%%%%%%%%%%%%%%%%%%%%%%%%%%%%%%%%%%%%%%%%%%%%%%%%%%%%%

Bisection is a bracketing method.

It maintains an interval known to contain a root.

Newton and secant methods are open methods.

They may converge much faster, but they do not guarantee that the root
remains bracketed.

Hybrid algorithms often combine both ideas.

%%%%%%%%%%%%%%%%%%%%%%%%%%%%%%%%%%%%%%%%%%%%%%%%%%%%%%%%%%%%
\subsection{Stopping Criteria for Root Finding}
\label{subsec:root-stopping-criteria}
%%%%%%%%%%%%%%%%%%%%%%%%%%%%%%%%%%%%%%%%%%%%%%%%%%%%%%%%%%%%

Possible stopping rules include:

\[
\boxed{
|x_{n+1}-x_n|
<
\varepsilon_x,
}
\]

\[
\boxed{
|f(x_n)|
<
\varepsilon_f,
}
\]

or a relative step test such as

\[
\boxed{
\frac{
|x_{n+1}-x_n|
}{
1+|x_{n+1}|
}
<
\varepsilon.
}
\]

A robust implementation often considers more than one criterion.

%%%%%%%%%%%%%%%%%%%%%%%%%%%%%%%%%%%%%%%%%%%%%%%%%%%%%%%%%%%%
\subsection{Numerical Differentiation}
\label{subsec:numerical-differentiation}
%%%%%%%%%%%%%%%%%%%%%%%%%%%%%%%%%%%%%%%%%%%%%%%%%%%%%%%%%%%%

The derivative is defined by

\[
f'(x)
=
\lim_{h\to0}
\frac{
f(x+h)-f(x)
}{
h
}.
\]

For finite \(h\),

\[
\boxed{
f'(x)
\approx
\frac{
f(x+h)-f(x)
}{
h
}.
}
\]

This is the forward-difference approximation.

%%%%%%%%%%%%%%%%%%%%%%%%%%%%%%%%%%%%%%%%%%%%%%%%%%%%%%%%%%%%
\subsection{Forward Difference}
\label{subsec:forward-difference}
%%%%%%%%%%%%%%%%%%%%%%%%%%%%%%%%%%%%%%%%%%%%%%%%%%%%%%%%%%%%

Taylor expansion gives

\[
f(x+h)
=
f(x)
+
hf'(x)
+
\frac{h^2}{2}f''(\xi).
\]

Therefore,

\[
\frac{
f(x+h)-f(x)
}{
h
}
=
f'(x)
+
O(h).
\]

Hence,

\[
\boxed{
f'(x)
=
\frac{
f(x+h)-f(x)
}{
h
}
+
O(h).
}
\]

The forward difference is first-order accurate.

%%%%%%%%%%%%%%%%%%%%%%%%%%%%%%%%%%%%%%%%%%%%%%%%%%%%%%%%%%%%
\subsection{Backward Difference}
\label{subsec:backward-difference}
%%%%%%%%%%%%%%%%%%%%%%%%%%%%%%%%%%%%%%%%%%%%%%%%%%%%%%%%%%%%

Similarly,

\[
\boxed{
f'(x)
=
\frac{
f(x)-f(x-h)
}{
h
}
+
O(h).
}
\]

This is the backward-difference approximation.

It is also first-order accurate.

%%%%%%%%%%%%%%%%%%%%%%%%%%%%%%%%%%%%%%%%%%%%%%%%%%%%%%%%%%%%
\subsection{Centered Difference}
\label{subsec:centered-difference}
%%%%%%%%%%%%%%%%%%%%%%%%%%%%%%%%%%%%%%%%%%%%%%%%%%%%%%%%%%%%

Using Taylor expansions for \(f(x+h)\) and \(f(x-h)\),

\[
f(x+h)
-
f(x-h)
=
2hf'(x)
+
O(h^3).
\]

Therefore,

\[
\boxed{
f'(x)
=
\frac{
f(x+h)-f(x-h)
}{
2h
}
+
O(h^2).
}
\]

The centered difference is second-order accurate.

%%%%%%%%%%%%%%%%%%%%%%%%%%%%%%%%%%%%%%%%%%%%%%%%%%%%%%%%%%%%
\subsection{Second Derivative Difference Formula}
\label{subsec:second-derivative-difference}
%%%%%%%%%%%%%%%%%%%%%%%%%%%%%%%%%%%%%%%%%%%%%%%%%%%%%%%%%%%%

Adding Taylor expansions yields

\[
f(x+h)
-
2f(x)
+
f(x-h)
=
h^2f''(x)
+
O(h^4).
\]

Thus,

\[
\boxed{
f''(x)
=
\frac{
f(x+h)
-
2f(x)
+
f(x-h)
}{
h^2
}
+
O(h^2).
}
\]

%%%%%%%%%%%%%%%%%%%%%%%%%%%%%%%%%%%%%%%%%%%%%%%%%%%%%%%%%%%%
\subsection{Worked Example: Centered Difference}
\label{subsec:worked-centered-difference}
%%%%%%%%%%%%%%%%%%%%%%%%%%%%%%%%%%%%%%%%%%%%%%%%%%%%%%%%%%%%

\begin{workedexamplebox}{Approximating a Derivative}

Let

\[
f(x)=e^x.
\]

Approximate

\[
f'(0)
\]

using centered difference with

\[
h=0.1.
\]

Then

\[
f'(0)
\approx
\frac{
e^{0.1}-e^{-0.1}
}{
0.2
}.
\]

Numerically,

\[
e^{0.1}\approx1.105170,
\]

\[
e^{-0.1}\approx0.904837.
\]

Therefore,

\[
f'(0)
\approx
\frac{
1.105170-0.904837
}{
0.2
}
\approx1.001665.
\]

The exact derivative is

\[
f'(0)=1.
\]

\begin{answerbox}
\[
\boxed{
f'(0)\approx1.001665.
}
\]
\end{answerbox}

\end{workedexamplebox}

%%%%%%%%%%%%%%%%%%%%%%%%%%%%%%%%%%%%%%%%%%%%%%%%%%%%%%%%%%%%
\subsection{Roundoff Versus Truncation in Differentiation}
\label{subsec:differentiation-roundoff-truncation}
%%%%%%%%%%%%%%%%%%%%%%%%%%%%%%%%%%%%%%%%%%%%%%%%%%%%%%%%%%%%

Decreasing \(h\) reduces truncation error.

However, very small \(h\) makes

\[
f(x+h)-f(x)
\]

a subtraction of nearly equal floating-point numbers.

This can amplify roundoff error.

Therefore,

\[
\boxed{
\text{smaller }h
\text{ is not always better}.
}
\]

An optimal step size balances truncation and roundoff effects.

%%%%%%%%%%%%%%%%%%%%%%%%%%%%%%%%%%%%%%%%%%%%%%%%%%%%%%%%%%%%
\subsection{Richardson Extrapolation}
\label{subsec:richardson-extrapolation}
%%%%%%%%%%%%%%%%%%%%%%%%%%%%%%%%%%%%%%%%%%%%%%%%%%%%%%%%%%%%

Suppose an approximation satisfies

\[
A(h)
=
A
+
Ch^p
+
O(h^{p+1}).
\]

Then

\[
A(h/2)
=
A
+
C
\frac{
h^p
}{
2^p
}
+
O(h^{p+1}).
\]

Eliminating the leading error term gives

\[
\boxed{
A_{\mathrm{rich}}
=
\frac{
2^pA(h/2)-A(h)
}{
2^p-1
}.
}
\]

This is Richardson extrapolation.

%%%%%%%%%%%%%%%%%%%%%%%%%%%%%%%%%%%%%%%%%%%%%%%%%%%%%%%%%%%%
\subsection{Numerical Integration}
\label{subsec:numerical-integration}
%%%%%%%%%%%%%%%%%%%%%%%%%%%%%%%%%%%%%%%%%%%%%%%%%%%%%%%%%%%%

Numerical integration approximates

\[
\boxed{
\int_a^b
f(x)\,dx.
}
\]

Quadrature methods replace \(f\) with simpler approximations whose
integrals can be computed exactly.

%%%%%%%%%%%%%%%%%%%%%%%%%%%%%%%%%%%%%%%%%%%%%%%%%%%%%%%%%%%%
\subsection{Midpoint Rule}
\label{subsec:midpoint-rule}
%%%%%%%%%%%%%%%%%%%%%%%%%%%%%%%%%%%%%%%%%%%%%%%%%%%%%%%%%%%%

The midpoint rule approximates

\[
\int_a^b
f(x)\,dx
\]

by using the function value at the midpoint

\[
m
=
\frac{a+b}{2}.
\]

Thus,

\[
\boxed{
\int_a^b
f(x)\,dx
\approx
(b-a)
f
\left(
\frac{a+b}{2}
\right).
}
\]

%%%%%%%%%%%%%%%%%%%%%%%%%%%%%%%%%%%%%%%%%%%%%%%%%%%%%%%%%%%%
\subsection{Trapezoidal Rule}
\label{subsec:trapezoidal-rule}
%%%%%%%%%%%%%%%%%%%%%%%%%%%%%%%%%%%%%%%%%%%%%%%%%%%%%%%%%%%%

Approximating \(f\) with the line connecting

\[
(a,f(a))
\]

and

\[
(b,f(b))
\]

gives

\[
\boxed{
\int_a^b
f(x)\,dx
\approx
\frac{
b-a
}{
2
}
[
f(a)+f(b)
].
}
\]

This is the trapezoidal rule.

%%%%%%%%%%%%%%%%%%%%%%%%%%%%%%%%%%%%%%%%%%%%%%%%%%%%%%%%%%%%
\subsection{Single-Interval Trapezoidal Error}
\label{subsec:trapezoidal-error-single}
%%%%%%%%%%%%%%%%%%%%%%%%%%%%%%%%%%%%%%%%%%%%%%%%%%%%%%%%%%%%

If \(f''\) is continuous, then for some \(\xi\in(a,b)\),

\[
\boxed{
E_T
=
-
\frac{
(b-a)^3
}{
12
}
f''(\xi).
}
\]

Thus the local error depends on curvature.

%%%%%%%%%%%%%%%%%%%%%%%%%%%%%%%%%%%%%%%%%%%%%%%%%%%%%%%%%%%%
\subsection{Simpson's Rule}
\label{subsec:simpsons-rule}
%%%%%%%%%%%%%%%%%%%%%%%%%%%%%%%%%%%%%%%%%%%%%%%%%%%%%%%%%%%%

Simpson's rule uses a quadratic interpolant through

\[
a,
\qquad
\frac{a+b}{2},
\qquad
b.
\]

The approximation is

\[
\boxed{
\int_a^b
f(x)\,dx
\approx
\frac{
b-a
}{
6
}
\left[
f(a)
+
4f
\left(
\frac{a+b}{2}
\right)
+
f(b)
\right].
}
\]

%%%%%%%%%%%%%%%%%%%%%%%%%%%%%%%%%%%%%%%%%%%%%%%%%%%%%%%%%%%%
\subsection{Worked Example: Simpson's Rule}
\label{subsec:worked-simpson}
%%%%%%%%%%%%%%%%%%%%%%%%%%%%%%%%%%%%%%%%%%%%%%%%%%%%%%%%%%%%

\begin{workedexamplebox}{Approximating \(\int_0^1 x^4\,dx\)}

The exact integral is

\[
\int_0^1
x^4\,dx
=
\frac15
=
0.2.
\]

Using Simpson's rule,

\[
S
=
\frac16
\left[
f(0)
+
4f(1/2)
+
f(1)
\right].
\]

Since

\[
f(0)=0,
\]

\[
f(1/2)=\frac1{16},
\]

and

\[
f(1)=1,
\]

we obtain

\[
S
=
\frac16
\left[
0
+
\frac14
+
1
\right]
=
\frac{5}{24}.
\]

Thus,

\[
S
\approx0.208333.
\]

\begin{answerbox}
\[
\boxed{
\int_0^1x^4\,dx
\approx0.208333
}
\]

using one Simpson panel.
\end{answerbox}

\end{workedexamplebox}

%%%%%%%%%%%%%%%%%%%%%%%%%%%%%%%%%%%%%%%%%%%%%%%%%%%%%%%%%%%%
\subsection{Composite Quadrature}
\label{subsec:composite-quadrature}
%%%%%%%%%%%%%%%%%%%%%%%%%%%%%%%%%%%%%%%%%%%%%%%%%%%%%%%%%%%%

Accuracy is improved by dividing

\[
[a,b]
\]

into smaller subintervals.

Let

\[
x_i
=
a+ih,
\]

where

\[
\boxed{
h
=
\frac{
b-a
}{
n
}.
}
\]

Apply the quadrature rule on each subinterval and sum.

%%%%%%%%%%%%%%%%%%%%%%%%%%%%%%%%%%%%%%%%%%%%%%%%%%%%%%%%%%%%
\subsection{Composite Trapezoidal Rule}
\label{subsec:composite-trapezoidal}
%%%%%%%%%%%%%%%%%%%%%%%%%%%%%%%%%%%%%%%%%%%%%%%%%%%%%%%%%%%%

The composite trapezoidal rule is

\[
\boxed{
T_n
=
h
\left[
\frac12f(x_0)
+
\sum_{i=1}^{n-1}
f(x_i)
+
\frac12f(x_n)
\right].
}
\]

For sufficiently smooth \(f\),

\[
\boxed{
E_T
=
O(h^2).
}
\]

%%%%%%%%%%%%%%%%%%%%%%%%%%%%%%%%%%%%%%%%%%%%%%%%%%%%%%%%%%%%
\subsection{Composite Midpoint Rule}
\label{subsec:composite-midpoint}
%%%%%%%%%%%%%%%%%%%%%%%%%%%%%%%%%%%%%%%%%%%%%%%%%%%%%%%%%%%%

If each subinterval has width \(h\), let \(m_i\) be its midpoint.

Then

\[
\boxed{
M_n
=
h
\sum_{i=1}^{n}
f(m_i).
}
\]

For sufficiently smooth \(f\),

\[
\boxed{
E_M
=
O(h^2).
}
\]

%%%%%%%%%%%%%%%%%%%%%%%%%%%%%%%%%%%%%%%%%%%%%%%%%%%%%%%%%%%%
\subsection{Composite Simpson's Rule}
\label{subsec:composite-simpson}
%%%%%%%%%%%%%%%%%%%%%%%%%%%%%%%%%%%%%%%%%%%%%%%%%%%%%%%%%%%%

Assume \(n\) is even.

Then

\[
\boxed{
S_n
=
\frac{h}{3}
\left[
f(x_0)
+
4
\sum_{\substack{i=1\\i\text{ odd}}}^{n-1}
f(x_i)
+
2
\sum_{\substack{i=2\\i\text{ even}}}^{n-2}
f(x_i)
+
f(x_n)
\right].
}
\]

For sufficiently smooth \(f\),

\[
\boxed{
E_S
=
O(h^4).
}
\]

%%%%%%%%%%%%%%%%%%%%%%%%%%%%%%%%%%%%%%%%%%%%%%%%%%%%%%%%%%%%
\subsection{Composite Simpson Error Bound}
\label{subsec:composite-simpson-error}
%%%%%%%%%%%%%%%%%%%%%%%%%%%%%%%%%%%%%%%%%%%%%%%%%%%%%%%%%%%%

If \(f^{(4)}\) is continuous,

\[
\boxed{
E_S
=
-
\frac{
b-a
}{
180
}
h^4
f^{(4)}(\xi)
}
\]

for some \(\xi\in(a,b)\).

Therefore,

\[
\boxed{
|E_S|
\leq
\frac{
b-a
}{
180
}
h^4
\max_{x\in[a,b]}
|f^{(4)}(x)|.
}
\]

%%%%%%%%%%%%%%%%%%%%%%%%%%%%%%%%%%%%%%%%%%%%%%%%%%%%%%%%%%%%
\subsection{Adaptive Quadrature}
\label{subsec:adaptive-quadrature}
%%%%%%%%%%%%%%%%%%%%%%%%%%%%%%%%%%%%%%%%%%%%%%%%%%%%%%%%%%%%

Uniform subdivision may waste computation in smooth regions.

Adaptive quadrature places more points where the integrand is difficult.

A typical strategy:

\begin{enumerate}
    \item approximate the integral on an interval;
    \item subdivide the interval;
    \item compare coarse and refined approximations;
    \item refine further if the estimated error is too large.
\end{enumerate}

%%%%%%%%%%%%%%%%%%%%%%%%%%%%%%%%%%%%%%%%%%%%%%%%%%%%%%%%%%%%
\subsection{Interpolation}
\label{subsec:numerical-interpolation}
%%%%%%%%%%%%%%%%%%%%%%%%%%%%%%%%%%%%%%%%%%%%%%%%%%%%%%%%%%%%

Suppose data are known at distinct nodes

\[
x_0,\ldots,x_n
\]

with values

\[
y_i=f(x_i).
\]

Polynomial interpolation seeks polynomial \(p_n\) of degree at most \(n\)
such that

\[
\boxed{
p_n(x_i)
=
y_i,
\qquad
i=0,\ldots,n.
}
\]

%%%%%%%%%%%%%%%%%%%%%%%%%%%%%%%%%%%%%%%%%%%%%%%%%%%%%%%%%%%%
\subsection{Existence and Uniqueness of Polynomial Interpolation}
\label{subsec:interpolation-existence-uniqueness}
%%%%%%%%%%%%%%%%%%%%%%%%%%%%%%%%%%%%%%%%%%%%%%%%%%%%%%%%%%%%

For \(n+1\) distinct nodes

\[
x_0,\ldots,x_n,
\]

there exists a unique polynomial of degree at most \(n\) interpolating the
corresponding data.

The uniqueness follows because two such polynomials would have a
difference with at least \(n+1\) distinct roots while having degree at most
\(n\).

%%%%%%%%%%%%%%%%%%%%%%%%%%%%%%%%%%%%%%%%%%%%%%%%%%%%%%%%%%%%
\subsection{Lagrange Basis Polynomials}
\label{subsec:lagrange-basis-polynomials}
%%%%%%%%%%%%%%%%%%%%%%%%%%%%%%%%%%%%%%%%%%%%%%%%%%%%%%%%%%%%

Define

\[
\boxed{
L_i(x)
=
\prod_{\substack{j=0\\j\neq i}}^{n}
\frac{
x-x_j
}{
x_i-x_j
}.
}
\]

Then

\[
\boxed{
L_i(x_j)
=
\begin{cases}
1,
&
i=j,
\\
0,
&
i\neq j.
\end{cases}
}
\]

These are the Lagrange basis polynomials.

%%%%%%%%%%%%%%%%%%%%%%%%%%%%%%%%%%%%%%%%%%%%%%%%%%%%%%%%%%%%
\subsection{Lagrange Interpolation Formula}
\label{subsec:lagrange-interpolation}
%%%%%%%%%%%%%%%%%%%%%%%%%%%%%%%%%%%%%%%%%%%%%%%%%%%%%%%%%%%%

The interpolation polynomial is

\[
\boxed{
p_n(x)
=
\sum_{i=0}^{n}
f(x_i)
L_i(x).
}
\]

The formula directly satisfies

\[
p_n(x_i)=f(x_i).
\]

%%%%%%%%%%%%%%%%%%%%%%%%%%%%%%%%%%%%%%%%%%%%%%%%%%%%%%%%%%%%
\subsection{Worked Example: Linear Interpolation}
\label{subsec:worked-linear-interpolation}
%%%%%%%%%%%%%%%%%%%%%%%%%%%%%%%%%%%%%%%%%%%%%%%%%%%%%%%%%%%%

\begin{workedexamplebox}{Interpolating Two Points}

Suppose

\[
f(1)=3,
\qquad
f(4)=9.
\]

The Lagrange basis functions are

\[
L_0(x)
=
\frac{x-4}{1-4},
\]

and

\[
L_1(x)
=
\frac{x-1}{4-1}.
\]

Therefore,

\[
p_1(x)
=
3
\frac{x-4}{-3}
+
9
\frac{x-1}{3}.
\]

Simplifying,

\[
p_1(x)
=
2x+1.
\]

Thus,

\begin{answerbox}
\[
\boxed{
p_1(x)=2x+1.
}
\]
\end{answerbox}

\end{workedexamplebox}

%%%%%%%%%%%%%%%%%%%%%%%%%%%%%%%%%%%%%%%%%%%%%%%%%%%%%%%%%%%%
\subsection{Interpolation Error Formula}
\label{subsec:interpolation-error-formula}
%%%%%%%%%%%%%%%%%%%%%%%%%%%%%%%%%%%%%%%%%%%%%%%%%%%%%%%%%%%%

Suppose \(f\) is sufficiently differentiable.

For each \(x\), there exists \(\xi\) in the relevant interval such that

\[
\boxed{
f(x)-p_n(x)
=
\frac{
f^{(n+1)}(\xi)
}{
(n+1)!
}
\prod_{i=0}^{n}
(x-x_i).
}
\]

The error depends on both:

\[
\boxed{
\text{smoothness of }f
}
\]

and

\[
\boxed{
\text{placement of interpolation nodes}.
}
\]

%%%%%%%%%%%%%%%%%%%%%%%%%%%%%%%%%%%%%%%%%%%%%%%%%%%%%%%%%%%%
\subsection{Newton Divided Differences}
\label{subsec:newton-divided-differences}
%%%%%%%%%%%%%%%%%%%%%%%%%%%%%%%%%%%%%%%%%%%%%%%%%%%%%%%%%%%%

The Newton interpolation form is

\[
\boxed{
p_n(x)
=
f[x_0]
+
f[x_0,x_1]
(x-x_0)
+
\cdots
}
\]

\[
\boxed{
\phantom{p_n(x)={}}
+
f[x_0,\ldots,x_n]
\prod_{j=0}^{n-1}
(x-x_j).
}
\]

The quantities

\[
f[x_i,\ldots,x_j]
\]

are divided differences.

%%%%%%%%%%%%%%%%%%%%%%%%%%%%%%%%%%%%%%%%%%%%%%%%%%%%%%%%%%%%
\subsection{First Divided Difference}
\label{subsec:first-divided-difference}
%%%%%%%%%%%%%%%%%%%%%%%%%%%%%%%%%%%%%%%%%%%%%%%%%%%%%%%%%%%%

For distinct nodes,

\[
\boxed{
f[x_i,x_j]
=
\frac{
f(x_j)-f(x_i)
}{
x_j-x_i
}.
}
\]

This is the slope of the secant line between the two data points.

%%%%%%%%%%%%%%%%%%%%%%%%%%%%%%%%%%%%%%%%%%%%%%%%%%%%%%%%%%%%
\subsection{Higher Divided Differences}
\label{subsec:higher-divided-differences}
%%%%%%%%%%%%%%%%%%%%%%%%%%%%%%%%%%%%%%%%%%%%%%%%%%%%%%%%%%%%

Higher divided differences are defined recursively:

\[
\boxed{
f[x_i,\ldots,x_{i+k}]
=
\frac{
f[x_{i+1},\ldots,x_{i+k}]
-
f[x_i,\ldots,x_{i+k-1}]
}{
x_{i+k}-x_i
}.
}
\]

This creates a triangular divided-difference table.

%%%%%%%%%%%%%%%%%%%%%%%%%%%%%%%%%%%%%%%%%%%%%%%%%%%%%%%%%%%%
\subsection{Advantage of Newton Form}
\label{subsec:newton-interpolation-advantage}
%%%%%%%%%%%%%%%%%%%%%%%%%%%%%%%%%%%%%%%%%%%%%%%%%%%%%%%%%%%%

If a new interpolation node is added, the existing Newton coefficients do
not need to be recomputed completely.

One new divided-difference term can be appended.

This makes Newton form convenient for incremental interpolation.

%%%%%%%%%%%%%%%%%%%%%%%%%%%%%%%%%%%%%%%%%%%%%%%%%%%%%%%%%%%%
\subsection{Runge's Phenomenon}
\label{subsec:runges-phenomenon}
%%%%%%%%%%%%%%%%%%%%%%%%%%%%%%%%%%%%%%%%%%%%%%%%%%%%%%%%%%%%

High-degree polynomial interpolation at equally spaced nodes can oscillate
strongly near interval endpoints.

This behavior is called Runge's phenomenon.

Increasing the polynomial degree does not necessarily improve the
approximation uniformly.

%%%%%%%%%%%%%%%%%%%%%%%%%%%%%%%%%%%%%%%%%%%%%%%%%%%%%%%%%%%%
\subsection{Chebyshev Nodes Preview}
\label{subsec:chebyshev-nodes}
%%%%%%%%%%%%%%%%%%%%%%%%%%%%%%%%%%%%%%%%%%%%%%%%%%%%%%%%%%%%

One way to reduce large interpolation oscillations is to use nodes clustered
near interval endpoints.

Chebyshev-type nodes on \([-1,1]\) can be written as

\[
\boxed{
x_k
=
\cos
\left(
\frac{
(2k+1)\pi
}{
2(n+1)
}
\right),
\qquad
k=0,\ldots,n.
}
\]

These nodes reduce the maximum growth of the interpolation-node product.

%%%%%%%%%%%%%%%%%%%%%%%%%%%%%%%%%%%%%%%%%%%%%%%%%%%%%%%%%%%%
\subsection{Piecewise Polynomial Interpolation}
\label{subsec:piecewise-polynomial-interpolation}
%%%%%%%%%%%%%%%%%%%%%%%%%%%%%%%%%%%%%%%%%%%%%%%%%%%%%%%%%%%%

Instead of fitting one high-degree polynomial over the entire interval,
one can fit low-degree polynomials over subintervals.

This gives:

\[
\boxed{
\text{better local control}
}
\]

and often

\[
\boxed{
\text{better numerical behavior}.
}
\]

Piecewise linear interpolation is the simplest example.

%%%%%%%%%%%%%%%%%%%%%%%%%%%%%%%%%%%%%%%%%%%%%%%%%%%%%%%%%%%%
\subsection{Splines}
\label{subsec:numerical-splines}
%%%%%%%%%%%%%%%%%%%%%%%%%%%%%%%%%%%%%%%%%%%%%%%%%%%%%%%%%%%%

A spline is a piecewise polynomial with smoothness conditions at internal
knots.

Cubic splines use polynomials of degree at most \(3\) on each interval.

A standard cubic spline is typically required to have continuous:

\[
\boxed{
S(x),
}
\]

\[
\boxed{
S'(x),
}
\]

and

\[
\boxed{
S''(x).
}
\]

%%%%%%%%%%%%%%%%%%%%%%%%%%%%%%%%%%%%%%%%%%%%%%%%%%%%%%%%%%%%
\subsection{Natural Cubic Spline}
\label{subsec:natural-cubic-spline}
%%%%%%%%%%%%%%%%%%%%%%%%%%%%%%%%%%%%%%%%%%%%%%%%%%%%%%%%%%%%

A natural cubic spline imposes boundary conditions

\[
\boxed{
S''(x_0)=0,
}
\]

and

\[
\boxed{
S''(x_n)=0.
}
\]

These conditions complete the system needed to determine the spline.

Other endpoint conditions are also possible.

%%%%%%%%%%%%%%%%%%%%%%%%%%%%%%%%%%%%%%%%%%%%%%%%%%%%%%%%%%%%
\subsection{Interpolation Versus Approximation}
\label{subsec:interpolation-versus-approximation}
%%%%%%%%%%%%%%%%%%%%%%%%%%%%%%%%%%%%%%%%%%%%%%%%%%%%%%%%%%%%

Interpolation requires

\[
\boxed{
p(x_i)=y_i
}
\]

at every data point.

Approximation does not.

For noisy data, passing exactly through every observation may be
undesirable.

Least-squares fitting instead seeks a function that minimizes aggregate
error.

%%%%%%%%%%%%%%%%%%%%%%%%%%%%%%%%%%%%%%%%%%%%%%%%%%%%%%%%%%%%
\subsection{Least-Squares Approximation}
\label{subsec:numerical-least-squares}
%%%%%%%%%%%%%%%%%%%%%%%%%%%%%%%%%%%%%%%%%%%%%%%%%%%%%%%%%%%%

Suppose a model

\[
p(x;\theta)
\]

depends on parameters \(\theta\).

Given data

\[
(x_i,y_i),
\]

least squares minimizes

\[
\boxed{
S(\theta)
=
\sum_{i=1}^{m}
[
y_i-p(x_i;\theta)
]^2.
}
\]

Unlike interpolation, the residuals generally do not vanish individually.

%%%%%%%%%%%%%%%%%%%%%%%%%%%%%%%%%%%%%%%%%%%%%%%%%%%%%%%%%%%%
\subsection{Linear Least Squares}
\label{subsec:linear-least-squares-review}
%%%%%%%%%%%%%%%%%%%%%%%%%%%%%%%%%%%%%%%%%%%%%%%%%%%%%%%%%%%%

For

\[
A\mathbf{x}
\approx
\mathbf{b},
\]

least squares minimizes

\[
\boxed{
\|A\mathbf{x}-\mathbf{b}\|_2^2.
}
\]

The normal equations are

\[
\boxed{
A^TA\mathbf{x}
=
A^T\mathbf{b}.
}
\]

However, numerical linear algebra often prefers QR or SVD-based methods
because explicitly forming \(A^TA\) can worsen conditioning.

%%%%%%%%%%%%%%%%%%%%%%%%%%%%%%%%%%%%%%%%%%%%%%%%%%%%%%%%%%%%
\subsection{Residual Vector}
\label{subsec:least-squares-residual-vector}
%%%%%%%%%%%%%%%%%%%%%%%%%%%%%%%%%%%%%%%%%%%%%%%%%%%%%%%%%%%%

The residual vector is

\[
\boxed{
r
=
\mathbf{b}
-
A\widehat{\mathbf{x}}.
}
\]

At a linear least-squares solution,

\[
\boxed{
A^Tr=0.
}
\]

Thus the residual is orthogonal to the column space of \(A\).

%%%%%%%%%%%%%%%%%%%%%%%%%%%%%%%%%%%%%%%%%%%%%%%%%%%%%%%%%%%%
\subsection{Consistency}
\label{subsec:numerical-consistency}
%%%%%%%%%%%%%%%%%%%%%%%%%%%%%%%%%%%%%%%%%%%%%%%%%%%%%%%%%%%%

A discretization is consistent if the discrete equations approach the
original mathematical equations as the discretization parameter tends to
zero.

Informally,

\[
\boxed{
\text{discrete model}
\to
\text{continuous model}
}
\]

as

\[
\boxed{
h\to0.
}
\]

Consistency alone does not guarantee a useful numerical solution.

%%%%%%%%%%%%%%%%%%%%%%%%%%%%%%%%%%%%%%%%%%%%%%%%%%%%%%%%%%%%
\subsection{Stability}
\label{subsec:numerical-method-stability}
%%%%%%%%%%%%%%%%%%%%%%%%%%%%%%%%%%%%%%%%%%%%%%%%%%%%%%%%%%%%

A numerical scheme is stable when perturbations such as rounding,
discretization, or initial errors do not grow uncontrollably through the
algorithm.

The precise definition depends on the numerical problem.

For differential-equation methods, stability often describes the
propagation of perturbations across time steps.

%%%%%%%%%%%%%%%%%%%%%%%%%%%%%%%%%%%%%%%%%%%%%%%%%%%%%%%%%%%%
\subsection{Convergence}
\label{subsec:numerical-method-convergence}
%%%%%%%%%%%%%%%%%%%%%%%%%%%%%%%%%%%%%%%%%%%%%%%%%%%%%%%%%%%%

A numerical method is convergent if the computed approximation approaches
the exact mathematical solution as the discretization is refined.

Symbolically,

\[
\boxed{
u_h
\to
u
\qquad
\text{as}
\qquad
h\to0.
}
\]

Convergence is the ultimate approximation property.

%%%%%%%%%%%%%%%%%%%%%%%%%%%%%%%%%%%%%%%%%%%%%%%%%%%%%%%%%%%%
\subsection{Consistency, Stability, and Convergence}
\label{subsec:consistency-stability-convergence}
%%%%%%%%%%%%%%%%%%%%%%%%%%%%%%%%%%%%%%%%%%%%%%%%%%%%%%%%%%%%

These concepts play distinct roles:

\[
\boxed{
\text{consistency}
\rightarrow
\text{correct limiting equations},
}
\]

\[
\boxed{
\text{stability}
\rightarrow
\text{controlled error propagation},
}
\]

and

\[
\boxed{
\text{convergence}
\rightarrow
\text{approach to the true solution}.
}
\]

Their exact relationship depends on the numerical setting.

For linear well-posed initial-value discretizations, important theorems
connect stability and consistency with convergence.

%%%%%%%%%%%%%%%%%%%%%%%%%%%%%%%%%%%%%%%%%%%%%%%%%%%%%%%%%%%%
\subsection{A Priori Error Analysis}
\label{subsec:a-priori-error-analysis}
%%%%%%%%%%%%%%%%%%%%%%%%%%%%%%%%%%%%%%%%%%%%%%%%%%%%%%%%%%%%

An a priori error estimate predicts error using mathematical properties
before the numerical solution is known.

For example,

\[
\boxed{
|E_h|
\leq
Ch^p.
}
\]

This helps choose a mesh size or step size before computation.

%%%%%%%%%%%%%%%%%%%%%%%%%%%%%%%%%%%%%%%%%%%%%%%%%%%%%%%%%%%%
\subsection{A Posteriori Error Analysis}
\label{subsec:a-posteriori-error-analysis}
%%%%%%%%%%%%%%%%%%%%%%%%%%%%%%%%%%%%%%%%%%%%%%%%%%%%%%%%%%%%

An a posteriori error estimate uses the computed solution itself.

Examples include:

\[
\boxed{
\text{residual estimates},
}
\]

\[
\boxed{
\text{comparison of coarse and fine meshes},
}
\]

or

\[
\boxed{
\text{embedded methods}.
}
\]

A posteriori estimates are important for adaptive algorithms.

%%%%%%%%%%%%%%%%%%%%%%%%%%%%%%%%%%%%%%%%%%%%%%%%%%%%%%%%%%%%
\subsection{Empirical Convergence Rate}
\label{subsec:empirical-convergence-rate}
%%%%%%%%%%%%%%%%%%%%%%%%%%%%%%%%%%%%%%%%%%%%%%%%%%%%%%%%%%%%

Suppose

\[
E(h)
\approx
Ch^p.
\]

Then

\[
E(h/2)
\approx
C
\left(
\frac{h}{2}
\right)^p.
\]

Hence,

\[
\frac{
E(h)
}{
E(h/2)
}
\approx
2^p.
\]

Therefore,

\[
\boxed{
p
\approx
\log_2
\left(
\frac{
E(h)
}{
E(h/2)
}
\right).
}
\]

This can estimate convergence order numerically.

%%%%%%%%%%%%%%%%%%%%%%%%%%%%%%%%%%%%%%%%%%%%%%%%%%%%%%%%%%%%
\subsection{Worked Example: Experimental Order}
\label{subsec:worked-experimental-order}
%%%%%%%%%%%%%%%%%%%%%%%%%%%%%%%%%%%%%%%%%%%%%%%%%%%%%%%%%%%%

\begin{workedexamplebox}{Estimating Convergence Rate}

Suppose the numerical errors are

\[
E(h)=0.008,
\]

and

\[
E(h/2)=0.002.
\]

Then

\[
\frac{
E(h)
}{
E(h/2)
}
=
4.
\]

Thus,

\[
p
\approx
\log_2(4)
=
2.
\]

\begin{answerbox}
\[
\boxed{
p\approx2.
}
\]
\end{answerbox}

The observed method is approximately second-order accurate over these mesh
sizes.

\end{workedexamplebox}

%%%%%%%%%%%%%%%%%%%%%%%%%%%%%%%%%%%%%%%%%%%%%%%%%%%%%%%%%%%%
\subsection{Algorithmic Complexity}
\label{subsec:numerical-algorithmic-complexity}
%%%%%%%%%%%%%%%%%%%%%%%%%%%%%%%%%%%%%%%%%%%%%%%%%%%%%%%%%%%%

Numerical methods must be evaluated not only by accuracy but also by
computational cost.

Possible cost measures include:

\[
\boxed{
\text{arithmetic operations},
}
\]

\[
\boxed{
\text{memory usage},
}
\]

\[
\boxed{
\text{iterations},
}
\]

and

\[
\boxed{
\text{communication cost}.
}
\]

A more accurate method is not always preferable if it is prohibitively
expensive.

%%%%%%%%%%%%%%%%%%%%%%%%%%%%%%%%%%%%%%%%%%%%%%%%%%%%%%%%%%%%
\subsection{Accuracy Versus Cost}
\label{subsec:accuracy-versus-cost}
%%%%%%%%%%%%%%%%%%%%%%%%%%%%%%%%%%%%%%%%%%%%%%%%%%%%%%%%%%%%

Suppose a method has error

\[
E(h)
=
O(h^p).
\]

A smaller \(h\) improves accuracy but usually increases the number of
computational points.

Thus numerical analysis seeks a balance between:

\[
\boxed{
\text{accuracy}
}
\]

and

\[
\boxed{
\text{computational effort}.
}
\]

%%%%%%%%%%%%%%%%%%%%%%%%%%%%%%%%%%%%%%%%%%%%%%%%%%%%%%%%%%%%
\subsection{Verification}
\label{subsec:numerical-verification}
%%%%%%%%%%%%%%%%%%%%%%%%%%%%%%%%%%%%%%%%%%%%%%%%%%%%%%%%%%%%

Verification asks:

\[
\boxed{
\text{Are we solving the mathematical model correctly?}
}
\]

Verification methods include:

\[
\boxed{
\text{mesh refinement},
}
\]

\[
\boxed{
\text{convergence studies},
}
\]

\[
\boxed{
\text{comparison with exact benchmark solutions},
}
\]

and

\[
\boxed{
\text{residual checks}.
}
\]

%%%%%%%%%%%%%%%%%%%%%%%%%%%%%%%%%%%%%%%%%%%%%%%%%%%%%%%%%%%%
\subsection{Validation}
\label{subsec:numerical-validation}
%%%%%%%%%%%%%%%%%%%%%%%%%%%%%%%%%%%%%%%%%%%%%%%%%%%%%%%%%%%%

Validation asks:

\[
\boxed{
\text{Is the mathematical model an adequate representation of the real
system?}
}
\]

A numerical algorithm may solve a model accurately even if the model
itself poorly represents reality.

Thus:

\[
\boxed{
\text{verification}
\neq
\text{validation}.
}
\]

%%%%%%%%%%%%%%%%%%%%%%%%%%%%%%%%%%%%%%%%%%%%%%%%%%%%%%%%%%%%
\subsection{Reproducibility}
\label{subsec:numerical-reproducibility}
%%%%%%%%%%%%%%%%%%%%%%%%%%%%%%%%%%%%%%%%%%%%%%%%%%%%%%%%%%%%

A numerical result should ideally document:

\[
\boxed{
\text{model equations},
}
\]

\[
\boxed{
\text{parameter values},
}
\]

\[
\boxed{
\text{initial and boundary conditions},
}
\]

\[
\boxed{
\text{algorithm},
}
\]

\[
\boxed{
\text{tolerances},
}
\]

and

\[
\boxed{
\text{software or precision assumptions}.
}
\]

This allows results to be independently checked.

%%%%%%%%%%%%%%%%%%%%%%%%%%%%%%%%%%%%%%%%%%%%%%%%%%%%%%%%%%%%
\subsection{Numerical Analysis Workflow}
\label{subsec:numerical-analysis-workflow}
%%%%%%%%%%%%%%%%%%%%%%%%%%%%%%%%%%%%%%%%%%%%%%%%%%%%%%%%%%%%

A practical workflow is:

\begin{enumerate}
    \item identify the mathematical problem;
    \item determine whether an exact solution is available;
    \item identify the desired numerical quantity;
    \item analyze conditioning;
    \item choose an appropriate numerical method;
    \item derive or understand its truncation error;
    \item select tolerances and discretization parameters;
    \item compute the approximation;
    \item monitor residuals and iteration progress;
    \item refine the discretization;
    \item estimate convergence order;
    \item compare independent methods when possible;
    \item check sensitivity to data and parameters;
    \item inspect possible roundoff or cancellation;
    \item verify the numerical solution;
    \item validate the underlying mathematical model when appropriate.
\end{enumerate}

%%%%%%%%%%%%%%%%%%%%%%%%%%%%%%%%%%%%%%%%%%%%%%%%%%%%%%%%%%%%
\subsection{Choosing a Numerical Method}
\label{subsec:choosing-numerical-method}
%%%%%%%%%%%%%%%%%%%%%%%%%%%%%%%%%%%%%%%%%%%%%%%%%%%%%%%%%%%%

For a continuous scalar root with a reliable sign-changing interval,
consider:

\[
\boxed{
\text{bisection}.
}
\]

For rapid local root convergence with a derivative, consider:

\[
\boxed{
\text{Newton's method}.
}
\]

Without derivatives, consider:

\[
\boxed{
\text{secant methods}.
}
\]

For derivative approximation, centered differences often improve on
one-sided first-order formulas.

For smooth numerical integration, consider:

\[
\boxed{
\text{composite quadrature or adaptive quadrature}.
}
\]

For exact fitting of smooth tabulated data, consider:

\[
\boxed{
\text{interpolation or splines}.
}
\]

For noisy or overdetermined data, consider:

\[
\boxed{
\text{least squares}.
}
\]

%%%%%%%%%%%%%%%%%%%%%%%%%%%%%%%%%%%%%%%%%%%%%%%%%%%%%%%%%%%%
\subsection{Common Errors}
\label{subsec:numerical-analysis-common-errors}
%%%%%%%%%%%%%%%%%%%%%%%%%%%%%%%%%%%%%%%%%%%%%%%%%%%%%%%%%%%%

\subsubsection{Treating a Displayed Decimal as Exact}

A printed number such as

\[
1.41421356
\]

may only be an approximation.

Displayed digits should not automatically be interpreted as exact data.

%%%%%%%%%%%%%%%%%%%%%%%%%%%%%%%%%%%%%%%%%%%%%%%%%%%%%%%%%%%%
\subsubsection{Confusing Absolute and Relative Error}

An absolute error of \(1\) is enormous if the true value is \(10^{-3}\),
but tiny if the true value is \(10^{12}\).

Scale matters.

%%%%%%%%%%%%%%%%%%%%%%%%%%%%%%%%%%%%%%%%%%%%%%%%%%%%%%%%%%%%
\subsubsection{Assuming More Decimal Places Means More Accuracy}

Software can display digits that are not mathematically reliable.

Accuracy depends on the computation, conditioning, and input data.

%%%%%%%%%%%%%%%%%%%%%%%%%%%%%%%%%%%%%%%%%%%%%%%%%%%%%%%%%%%%
\subsubsection{Ignoring Floating-Point Representation}

Real-number algebra and floating-point arithmetic are not identical.

Operations such as

\[
(a+b)+c
\]

and

\[
a+(b+c)
\]

can produce slightly different floating-point results.

%%%%%%%%%%%%%%%%%%%%%%%%%%%%%%%%%%%%%%%%%%%%%%%%%%%%%%%%%%%%
\subsubsection{Subtracting Nearly Equal Numbers Unnecessarily}

This can cause catastrophic cancellation.

Algebraic reformulation may produce a more stable expression.

%%%%%%%%%%%%%%%%%%%%%%%%%%%%%%%%%%%%%%%%%%%%%%%%%%%%%%%%%%%%
\subsubsection{Confusing Conditioning and Stability}

An excellent algorithm cannot completely eliminate sensitivity inherent in
an ill-conditioned problem.

Likewise, a well-conditioned problem can still be solved badly by an
unstable algorithm.

%%%%%%%%%%%%%%%%%%%%%%%%%%%%%%%%%%%%%%%%%%%%%%%%%%%%%%%%%%%%
\subsubsection{Assuming Small Residual Means Small Error}

A small residual may correspond to a large forward error when the problem
is ill conditioned.

Residual and conditioning should be interpreted together.

%%%%%%%%%%%%%%%%%%%%%%%%%%%%%%%%%%%%%%%%%%%%%%%%%%%%%%%%%%%%
\subsubsection{Using Bisection Without a Valid Bracket}

Classical bisection requires a continuous function and a sign-changing
interval:

\[
\boxed{
f(a)f(b)<0.
}
\]

Without these conditions, the usual root guarantee is lost.

%%%%%%%%%%%%%%%%%%%%%%%%%%%%%%%%%%%%%%%%%%%%%%%%%%%%%%%%%%%%
\subsubsection{Using Newton's Method Without Checking the Derivative}

If

\[
f'(x_n)
\approx0,
\]

the Newton step can become extremely large or undefined.

%%%%%%%%%%%%%%%%%%%%%%%%%%%%%%%%%%%%%%%%%%%%%%%%%%%%%%%%%%%%
\subsubsection{Assuming Newton's Method Always Converges}

Newton's method is locally powerful but not globally guaranteed.

Initial guesses matter.

%%%%%%%%%%%%%%%%%%%%%%%%%%%%%%%%%%%%%%%%%%%%%%%%%%%%%%%%%%%%
\subsubsection{Using an Extremely Small Finite-Difference Step}

Truncation error decreases with \(h\), but roundoff and cancellation can
increase.

The best step is not always the smallest representable step.

%%%%%%%%%%%%%%%%%%%%%%%%%%%%%%%%%%%%%%%%%%%%%%%%%%%%%%%%%%%%
\subsubsection{Using Simpson's Rule With an Odd Number of Composite Subintervals}

The standard composite Simpson rule requires

\[
\boxed{
n
\text{ even}.
}
\]

%%%%%%%%%%%%%%%%%%%%%%%%%%%%%%%%%%%%%%%%%%%%%%%%%%%%%%%%%%%%
\subsubsection{Using One High-Degree Polynomial Automatically}

High-degree interpolation at equally spaced points can exhibit severe
oscillations.

Splines or better node placement may be preferable.

%%%%%%%%%%%%%%%%%%%%%%%%%%%%%%%%%%%%%%%%%%%%%%%%%%%%%%%%%%%%
\subsubsection{Confusing Interpolation and Regression}

Interpolation passes exactly through the supplied data.

Least-squares regression generally does not.

%%%%%%%%%%%%%%%%%%%%%%%%%%%%%%%%%%%%%%%%%%%%%%%%%%%%%%%%%%%%
\subsubsection{Ignoring Convergence Studies}

A numerical answer at one grid size or tolerance gives limited evidence of
accuracy.

Refinement studies help confirm convergence.

%%%%%%%%%%%%%%%%%%%%%%%%%%%%%%%%%%%%%%%%%%%%%%%%%%%%%%%%%%%%
\subsubsection{Confusing Verification and Validation}

A numerically accurate solution to the wrong model is still the wrong
modeling result.

Both questions should be considered separately.

%%%%%%%%%%%%%%%%%%%%%%%%%%%%%%%%%%%%%%%%%%%%%%%%%%%%%%%%%%%%
\subsection{Applications}
\label{subsec:numerical-analysis-applications}
%%%%%%%%%%%%%%%%%%%%%%%%%%%%%%%%%%%%%%%%%%%%%%%%%%%%%%%%%%%%

\begin{applicationbox}

\textbf{Engineering.}
Numerical methods solve nonlinear design equations, structural models,
fluid models, and control problems.

\medskip

\textbf{Physics.}
Many physical systems require numerical solutions of differential and
integral equations.

\medskip

\textbf{Statistics.}
Likelihood maximization, regression, Bayesian computation, and simulation
depend heavily on numerical algorithms.

\medskip

\textbf{Optimization.}
Modern optimization relies on numerical linear algebra, derivative
approximations, line searches, and iterative solution methods.

\medskip

\textbf{Differential Equations.}
ODE and PDE solutions are frequently available only through numerical
discretization.

\medskip

\textbf{Finance.}
Pricing, risk estimation, portfolio optimization, and stochastic models
require numerical integration, simulation, and equation solving.

\medskip

\textbf{Data Science.}
Large-scale fitting and machine-learning algorithms are fundamentally
numerical optimization procedures.

\medskip

\textbf{Scientific Computing.}
Numerical analysis provides the mathematical foundation for reliable
computational science.

\end{applicationbox}

%%%%%%%%%%%%%%%%%%%%%%%%%%%%%%%%%%%%%%%%%%%%%%%%%%%%%%%%%%%%
\subsection{Exercises}
\label{subsec:numerical-analysis-exercises}
%%%%%%%%%%%%%%%%%%%%%%%%%%%%%%%%%%%%%%%%%%%%%%%%%%%%%%%%%%%%

\begin{exercise}[1]
Define absolute error.
\end{exercise}

\begin{exercise}[1]
Define relative error.
\end{exercise}

\begin{exercise}[1]
Define roundoff error.
\end{exercise}

\begin{exercise}[1]
Define truncation error.
\end{exercise}

\begin{exercise}[1]
Define conditioning.
\end{exercise}

\begin{exercise}[1]
Define numerical stability.
\end{exercise}

\begin{exercise}[1]
Define forward error.
\end{exercise}

\begin{exercise}[1]
Define backward error.
\end{exercise}

\begin{exercise}[1]
Define a residual.
\end{exercise}

\begin{exercise}[1]
Define convergence order.
\end{exercise}

\begin{exercise}[2]
For exact value

\[
x=2.718281828
\]

and approximation

\[
\widehat{x}=2.718,
\]

compute absolute and relative error.
\end{exercise}

\begin{exercise}[2]
Explain the difference between overflow and underflow.
\end{exercise}

\begin{exercise}[2]
Give an example of catastrophic cancellation.
\end{exercise}

\begin{exercise}[2]
Use algebraic reformulation to avoid cancellation in

\[
\sqrt{x+1}-\sqrt{x}.
\]
\end{exercise}

\begin{exercise}[2]
Compute the relative condition number of

\[
f(x)=x^5.
\]
\end{exercise}

\begin{exercise}[3]
Compute the relative condition number of

\[
f(x)=\ln x.
\]
\end{exercise}

\begin{exercise}[3]
Compare forward error and residual for an approximate root.
\end{exercise}

\begin{exercise}[2]
Apply one iteration of bisection to a given sign-changing interval.
\end{exercise}

\begin{exercise}[3]
Use bisection to approximate a root of

\[
x^3-x-2=0
\]

to a specified tolerance.
\end{exercise}

\begin{exercise}[3]
Determine the number of bisection iterations needed to guarantee a
specified interval width.
\end{exercise}

\begin{exercise}[2]
Rewrite

\[
x^2-3=0
\]

as a fixed-point problem.
\end{exercise}

\begin{exercise}[3]
Determine whether a given fixed-point iteration is locally convergent by
examining \(g'(x^*)\).
\end{exercise}

\begin{exercise}[2]
Write Newton's iteration for

\[
f(x)=x^3-5.
\]
\end{exercise}

\begin{exercise}[3]
Use Newton's method to approximate

\[
\sqrt[3]{5}.
\]
\end{exercise}

\begin{exercise}[3]
Show that Newton's method for

\[
x^2-a=0
\]

gives

\[
x_{n+1}
=
\frac12
\left(
x_n+\frac{a}{x_n}
\right).
\]
\end{exercise}

\begin{exercise}[3]
Apply the secant method to a nonlinear equation using two initial points.
\end{exercise}

\begin{exercise}[3]
Compare the iteration counts of bisection, Newton, and secant methods on
the same root problem.
\end{exercise}

\begin{exercise}[3]
Construct a stopping rule using both step size and residual.
\end{exercise}

\begin{exercise}[2]
Derive the forward-difference approximation using Taylor expansion.
\end{exercise}

\begin{exercise}[2]
Derive the backward-difference approximation.
\end{exercise}

\begin{exercise}[3]
Derive the centered-difference formula.
\end{exercise}

\begin{exercise}[3]
Derive the centered second-derivative formula.
\end{exercise}

\begin{exercise}[3]
Use forward and centered differences to approximate

\[
\frac{d}{dx}\sin x
\]

at a specified point.
\end{exercise}

\begin{exercise}[3]
Compare their errors for several values of \(h\).
\end{exercise}

\begin{exercise}[3]
Explain why the error eventually increases when \(h\) becomes extremely
small in floating-point arithmetic.
\end{exercise}

\begin{exercise}[3]
Apply Richardson extrapolation to two derivative approximations.
\end{exercise}

\begin{exercise}[2]
Apply the midpoint rule to a definite integral.
\end{exercise}

\begin{exercise}[2]
Apply the trapezoidal rule to the same integral.
\end{exercise}

\begin{exercise}[2]
Apply Simpson's rule to the same integral.
\end{exercise}

\begin{exercise}[3]
Construct the composite trapezoidal rule with four subintervals.
\end{exercise}

\begin{exercise}[3]
Construct the composite Simpson rule with four subintervals.
\end{exercise}

\begin{exercise}[3]
Use a derivative bound to estimate composite trapezoidal error.
\end{exercise}

\begin{exercise}[3]
Use a fourth-derivative bound to estimate composite Simpson error.
\end{exercise}

\begin{exercise}[3]
Explain how adaptive quadrature decides where to refine.
\end{exercise}

\begin{exercise}[2]
Construct the linear interpolating polynomial through two points.
\end{exercise}

\begin{exercise}[3]
Construct a quadratic Lagrange interpolant through three points.
\end{exercise}

\begin{exercise}[3]
Verify that each Lagrange basis polynomial has the required nodal values.
\end{exercise}

\begin{exercise}[3]
Construct a divided-difference table.
\end{exercise}

\begin{exercise}[3]
Write the corresponding Newton interpolation polynomial.
\end{exercise}

\begin{exercise}[3]
Use the interpolation-error formula to bound the error at a specified
point.
\end{exercise}

\begin{exercise}[3]
Explain Runge's phenomenon.
\end{exercise}

\begin{exercise}[3]
Compare equally spaced nodes with Chebyshev-type nodes conceptually.
\end{exercise}

\begin{exercise}[3]
Construct a piecewise linear interpolant.
\end{exercise}

\begin{exercise}[3]
State the continuity requirements of a cubic spline.
\end{exercise}

\begin{exercise}[3]
Explain the natural cubic-spline boundary conditions.
\end{exercise}

\begin{exercise}[3]
Explain the difference between interpolation and least squares.
\end{exercise}

\begin{exercise}[3]
Write the normal equations for a linear least-squares problem.
\end{exercise}

\begin{exercise}[3]
Explain why QR factorization can be numerically preferable to normal
equations.
\end{exercise}

\begin{exercise}[3]
Estimate an empirical convergence rate from errors at \(h\) and \(h/2\).
\end{exercise}

\begin{exercise}[3]
Distinguish a priori and a posteriori error estimates.
\end{exercise}

\begin{exercise}[3]
Design a simple numerical convergence study.
\end{exercise}

\begin{exercise}[3]
Explain the difference between verification and validation.
\end{exercise}

\begin{exercise}[4]
Derive a rigorous error bound for bisection.
\end{exercise}

\begin{exercise}[4]
Prove local quadratic convergence of Newton's method for a simple root.
\end{exercise}

\begin{exercise}[4]
Study Newton convergence for multiple roots.
\end{exercise}

\begin{exercise}[4]
Derive the secant-method convergence order.
\end{exercise}

\begin{exercise}[4]
Study safeguarded Newton methods.
\end{exercise}

\begin{exercise}[4]
Study Brent-type hybrid root solvers.
\end{exercise}

\begin{exercise}[4]
Derive optimal finite-difference step-size scaling under roundoff error.
\end{exercise}

\begin{exercise}[4]
Derive high-order finite-difference formulas using Taylor expansions.
\end{exercise}

\begin{exercise}[4]
Study Richardson extrapolation systematically.
\end{exercise}

\begin{exercise}[4]
Derive Romberg integration from trapezoidal extrapolation.
\end{exercise}

\begin{exercise}[4]
Study Gaussian quadrature.
\end{exercise}

\begin{exercise}[4]
Derive Gauss--Legendre quadrature nodes and weights for low orders.
\end{exercise}

\begin{exercise}[4]
Study adaptive quadrature error estimation.
\end{exercise}

\begin{exercise}[4]
Prove uniqueness of polynomial interpolation.
\end{exercise}

\begin{exercise}[4]
Derive the interpolation-error formula.
\end{exercise}

\begin{exercise}[4]
Study Lebesgue constants and interpolation conditioning.
\end{exercise}

\begin{exercise}[4]
Study Chebyshev interpolation and near-minimax approximation.
\end{exercise}

\begin{exercise}[4]
Derive cubic-spline equations.
\end{exercise}

\begin{exercise}[4]
Study best approximation in normed spaces.
\end{exercise}

\begin{exercise}[4]
Study backward stability of numerical algorithms.
\end{exercise}

\begin{exercise}[4]
Investigate mixed forward-backward error analysis.
\end{exercise}

\begin{exercise}[4]
Study interval arithmetic for verified numerical computation.
\end{exercise}

%%%%%%%%%%%%%%%%%%%%%%%%%%%%%%%%%%%%%%%%%%%%%%%%%%%%%%%%%%%%
\subsection{Conceptual Summary}
\label{subsec:numerical-analysis-summary}
%%%%%%%%%%%%%%%%%%%%%%%%%%%%%%%%%%%%%%%%%%%%%%%%%%%%%%%%%%%%

Numerical analysis studies reliable approximation of mathematical
problems.

Approximation error may be measured using

\[
\boxed{
E_{\mathrm{abs}}
=
|x-\widehat{x}|
}
\]

and

\[
\boxed{
E_{\mathrm{rel}}
=
\frac{
|x-\widehat{x}|
}{
|x|
}.
}
\]

Important error sources include

\[
\boxed{
\text{roundoff}
+
\text{truncation}
+
\text{data error}.
}
\]

A problem's sensitivity is described by

\[
\boxed{
\text{conditioning},
}
\]

while an algorithm's error behavior is described by

\[
\boxed{
\text{stability}.
}
\]

Root-finding methods include

\[
\boxed{
\text{bisection},
}
\]

\[
\boxed{
\text{Newton},
}
\]

and

\[
\boxed{
\text{secant}.
}
\]

Numerical differentiation uses formulas such as

\[
\boxed{
f'(x)
\approx
\frac{
f(x+h)-f(x-h)
}{
2h
}.
}
\]

Numerical integration includes

\[
\boxed{
\text{midpoint},
}
\]

\[
\boxed{
\text{trapezoidal},
}
\]

and

\[
\boxed{
\text{Simpson}
}
\]

rules.

Polynomial interpolation can be written as

\[
\boxed{
p_n(x)
=
\sum_{i=0}^{n}
f(x_i)L_i(x).
}
\]

Reliable numerical computation requires consideration of

\[
\boxed{
\text{accuracy}
+
\text{conditioning}
+
\text{stability}
+
\text{convergence}
+
\text{computational cost}.
}
\]

%%%%%%%%%%%%%%%%%%%%%%%%%%%%%%%%%%%%%%%%%%%%%%%%%%%%%%%%%%%%
\subsection{Section Connections}
\label{subsec:numerical-analysis-connections}
%%%%%%%%%%%%%%%%%%%%%%%%%%%%%%%%%%%%%%%%%%%%%%%%%%%%%%%%%%%%

\chapterconnection{
Numerical Analysis begins Chapter~11 by establishing the mathematical
principles required for reliable computation: error measurement,
floating-point arithmetic, conditioning, stability, convergence,
root finding, numerical differentiation, quadrature, and interpolation.
These concepts underlie nearly every later computational topic. The next
section, Numerical Linear Algebra, develops stable and efficient methods
for solving linear systems, matrix factorizations, least-squares problems,
eigenvalue problems, and large sparse matrix computations.
}

%%%%%%%%%%%%%%%%%%%%%%%%%%%%%%%%%%%%%%%%%%%%%%%%%%%%%%%%%%%%
% SECTION SOURCE TRACKING
%%%%%%%%%%%%%%%%%%%%%%%%%%%%%%%%%%%%%%%%%%%%%%%%%%%%%%%%%%%%

%------------------------------------------------------------
% 11.1 Numerical Analysis
%------------------------------------------------------------
%
% Source basis:
%   - Standard introductory and advanced numerical-analysis theory.
%   - Standard floating-point, approximation, root-finding,
%     differentiation, quadrature, and interpolation methods.
%
% Used for:
%   - numerical approximation;
%   - absolute and relative error;
%   - significant digits;
%   - floating-point representation;
%   - finite precision;
%   - machine epsilon and unit roundoff;
%   - overflow and underflow;
%   - roundoff and truncation errors;
%   - catastrophic cancellation;
%   - Taylor-series analysis;
%   - Big-O error notation;
%   - conditioning;
%   - forward and backward error;
%   - numerical stability;
%   - residuals;
%   - bisection;
%   - fixed-point iteration;
%   - Newton's method;
%   - secant method;
%   - convergence order;
%   - root-finding stopping criteria;
%   - finite differences;
%   - forward, backward, and centered derivatives;
%   - second-derivative approximation;
%   - Richardson extrapolation;
%   - midpoint, trapezoidal, and Simpson quadrature;
%   - composite quadrature;
%   - adaptive quadrature;
%   - polynomial interpolation;
%   - Lagrange interpolation;
%   - Newton divided differences;
%   - interpolation error;
%   - Runge's phenomenon;
%   - Chebyshev-node preview;
%   - splines;
%   - least-squares approximation;
%   - consistency, stability, and convergence;
%   - a priori and a posteriori error estimates;
%   - empirical convergence rates;
%   - algorithmic complexity;
%   - verification and validation;
%   - reproducibility;
%   - applications and exercises.
%
% Notes:
%   - Numerical Linear Algebra is developed next in Section 11.2.
%   - Numerical Optimization follows in Section 11.3.
%   - Numerical ODEs and PDEs are developed in Sections 11.4 and 11.5.
%   - Monte Carlo Methods are developed in Section 11.6.
%
%%%%%%%%%%%%%%%%%%%%%%%%%%%%%%%%%%%%%%%%%%%%%%%%%%%%%%%%%%%%